\section{Related Work}
\label{sec:related}

We review related work from five perspectives: graph neural networks for graph classification, message passing architectures and operators, deep GNN training challenges, multi-operator and hybrid architectures, and cross-layer aggregation mechanisms. Our work builds upon these foundations by introducing a unified framework that integrates multiple graph operators through cross-branch residual connections.

\subsection{Graph Neural Networks for Graph Classification}

Graph-level prediction requires aggregating node-level features into fixed-size graph representations. Early graph neural networks \cite{gori2005new} laid the foundation for learning on graph-structured data through recursive message passing. The Message Passing Neural Networks (MPNN) framework \cite{gilmer2017neural} unified various GNN architectures under a common formalism, identifying message and vertex update functions as core components.

For graph classification, readout functions play a crucial role in converting node embeddings to graph-level representations. Simple approaches apply global pooling operations such as mean or sum pooling \cite{duvenaud2015convolutional}. However, these methods may not adequately capture hierarchical structural information. To address this, differentiable pooling methods have been proposed: DiffPool \cite{ying2018hierarchical} learns node cluster assignments to coarsen graphs hierarchically, SAGPool \cite{lee2019self} employs self-attention to identify salient nodes, and LocalPool \cite{cangea2018towards} enables sparse hierarchical classification. Recent surveys provide comprehensive overviews of graph pooling techniques \cite{liu2022graph,li2024graph}. While these approaches achieve strong performance, they primarily focus on intra-branch architectural designs and do not explicitly consider how multiple operators with different inductive biases can be integrated and coordinated through cross-branch information exchange.

\subsection{Message Passing Architectures and Operators}

Various message passing operators have been developed, each imposing different inductive biases on neighborhood aggregation. Graph Convolutional Networks (GCN) \cite{kipf2016semi} approximate spectral convolutions through normalized adjacency matrices, performing low-pass filtering that smooths node features. GraphSAGE \cite{hamilton2017inductive} samples and aggregates features from local neighborhoods using mean, LSTM, or pooling aggregators, enabling inductive learning on unseen graphs. Graph Isomorphism Networks (GIN) \cite{xu2018powerful} use injective aggregation functions to achieve maximum discriminative power, matching the expressive power of the Weisfeiler-Lehman test.

Attention mechanisms have been widely incorporated to weight neighbor contributions adaptively. Graph Attention Networks (GAT) \cite{velivckovic2017graph} compute attention coefficients for each neighbor, allowing nodes to focus on relevant information. Personalized Graph Neural Networks (PGNN) \cite{zhang2019personalized} further extend attention mechanisms for personalized aggregation. Recent work generalizes Transformer architectures to graphs \cite{dwivedi2020generalization} and incorporates edge directionality \cite{beaini2020directional, rossi2023edge} for improved modeling of directed and heterophilic graphs.

Other operator variants include ARMA filters \cite{bianchi2020graph} for flexible frequency response, MixHop \cite{abu2019mixhop} for mixing multi-hop neighborhoods, and anti-symmetric convolutions \cite{gravina2022anti} for stable deep architectures. While these operators are typically studied and deployed individually within homogeneous architectures, a critical gap remains: existing frameworks treat operators as alternative choices rather than complementary components that can be systematically combined and coordinated to leverage their diverse inductive biases within a unified model.

\subsection{Deep GNN Training and Representation Degradation}

Training deep GNNs faces fundamental challenges due to representation degradation. The \textit{over-smoothing} phenomenon causes node representations to converge toward indistinguishable values as depth increases \cite{li2018deeper, chen2023revisiting, li2023non}. Theoretical analysis \cite{li2023non} provides non-asymptotic bounds on over-smoothing, revealing fundamental limitations in expressive power for deep architectures. Additionally, \textit{over-squashing} \cite{topping2021understanding, alon2020bottleneck} occurs when long-range dependencies cannot be effectively propagated due to graph bottlenecks, further constraining model performance.

To enable deep training, various techniques have been proposed. Residual connections \cite{he2016deep}, originally developed for computer vision, have been adapted to GNNs through residual gated graph convolutions \cite{bresson2017residual}. Normalization methods such as PairNorm \cite{zhao2020pair} and GraphNorm \cite{cai2021graphnorm} alleviate over-smoothing by normalizing node features, while learnable normalization \cite{li2022learning} adapts to different graph distributions. Regularization techniques like DropEdge \cite{rong2019dropedge} randomly drop edges during training to prevent overfitting and improve generalization.

Advanced architectures attempt to train extremely deep networks: techniques for training 1000-layer GNNs \cite{li2021training} and simplified deep GCN approaches \cite{oono2020simple} demonstrate that depth can improve performance with proper design. However, these methods primarily operate within single-branch architectures and do not explicitly model how information can be exchanged across multiple operator-specific branches to preserve multi-scale features and mitigate degradation at depth.

\subsection{Multi-Operator and Hybrid Architectures}

Combining multiple aggregation schemes or architectural components has shown promise in enhancing representational capacity. Multi-hop attention \cite{wang2021multi} enables information flow across longer distances, while ARMA filters \cite{bianchi2020graph} provide flexible frequency responses beyond polynomial filters. Hybrid architectures incorporate edge directionality \cite{beaini2020directional, rossi2023edge} to capture directed relationships, and anti-symmetric convolutions \cite{gravina2022anti} improve stability in deep networks.

Surveys on attention-based GNNs \cite{sun2023attention} and graph embeddings \cite{cai2020understanding} provide comprehensive taxonomies of architectural variants. DenseGNN \cite{du2024densegnn} connects all layers densely, similar to DenseNet in computer vision, to maintain information flow. Hierarchical pooling methods \cite{pham2021hierarchical, ying2018hierarchical} create multi-scale representations through graph coarsening. However, most existing multi-operator approaches treat different components as modular plug-ins that are combined sequentially or in parallel, without explicit mechanisms for structured information exchange between them during training. Our work addresses this gap by introducing cross-branch residual connections that enable operators to continuously exchange and refine information throughout the forward pass.

\subsection{Cross-Layer and Multi-Scale Aggregation}

Cross-layer connections aim to capture multi-scale information by aggregating features from different depths. Jumping Knowledge (JK) networks \cite{xu2018representation} adaptively select which layer's representation to use for each node, enabling flexible neighborhood range selection. Graph U-Net \cite{gao2019graph} adopts encoder-decoder architecture with skip connections between corresponding layers.

Propagation-based methods improve information flow across multiple hops: APPNP \cite{klicpera2018combining} combines personalized PageRank with neural predictions, while PairNorm \cite{zhao2020pair} and DropEdge \cite{rong2019dropedge} mitigate over-smoothing through normalization and regularization. Recent work on over-smoothing \cite{chen2023revisiting, zhao2020tackling} provides comprehensive analysis and solutions.

Despite these advances, most cross-layer methods operate within a single homogeneous branch using one type of operator. A fundamental limitation is that they do not consider how multiple branches with different operators can leverage complementary inductive biases through structured cross-branch communication. Our framework introduces cross-branch residual connections that enable node-level and graph-level information exchange between operator-specific branches, facilitating synergistic collaboration where, for example, the smoothing properties of GCN can stabilize the attention weights of GAT, while the selectivity of GAT can refine the low-pass features of GCN. This cross-branch coordination, combined with multi-scale aggregation within each branch, provides a more comprehensive approach to integrating diverse architectural biases for improved graph classification.
